% Options for packages loaded elsewhere
\PassOptionsToPackage{unicode}{hyperref}
\PassOptionsToPackage{hyphens}{url}
\PassOptionsToPackage{dvipsnames,svgnames,x11names}{xcolor}
\documentclass[
  10pt,
  fontset=fandol]{ctexart}
\usepackage{xcolor}
\usepackage[margin=16mm]{geometry}
\usepackage{amsmath,amssymb}
\setcounter{secnumdepth}{-\maxdimen} % remove section numbering
\usepackage{iftex}
\ifPDFTeX
  \usepackage[T1]{fontenc}
  \usepackage[utf8]{inputenc}
  \usepackage{textcomp} % provide euro and other symbols
\else % if luatex or xetex
  \usepackage{unicode-math} % this also loads fontspec
  \defaultfontfeatures{Scale=MatchLowercase}
  \defaultfontfeatures[\rmfamily]{Ligatures=TeX,Scale=1}
\fi
\usepackage{lmodern}
\ifPDFTeX\else
  % xetex/luatex font selection
\fi
% Use upquote if available, for straight quotes in verbatim environments
\IfFileExists{upquote.sty}{\usepackage{upquote}}{}
\IfFileExists{microtype.sty}{% use microtype if available
  \usepackage[]{microtype}
  \UseMicrotypeSet[protrusion]{basicmath} % disable protrusion for tt fonts
}{}
\makeatletter
\@ifundefined{KOMAClassName}{% if non-KOMA class
  \IfFileExists{parskip.sty}{%
    \usepackage{parskip}
  }{% else
    \setlength{\parindent}{0pt}
    \setlength{\parskip}{6pt plus 2pt minus 1pt}}
}{% if KOMA class
  \KOMAoptions{parskip=half}}
\makeatother
\setlength{\emergencystretch}{3em} % prevent overfull lines
\providecommand{\tightlist}{%
  \setlength{\itemsep}{0pt}\setlength{\parskip}{0pt}}
\usepackage{amsmath,amssymb,mathtools,needspace}
\xeCJKDeclareCharClass{CJK}{"2032,"0400->"04FF,"2160->"216B,"2460->"2473,"25A0,"25C7,"25CB,"25CF}
\IfFontExistsTF{Arial Unicode MS}{\xeCJKsetup{AutoFallBack=true}\setCJKfallbackfamilyfont{\CJKrmdefault}{Arial Unicode MS}}{}
\setcounter{secnumdepth}{0}
\setlength{\emergencystretch}{2em}
\allowdisplaybreaks
\usepackage{bookmark}
\IfFileExists{xurl.sty}{\usepackage{xurl}}{} % add URL line breaks if available
\urlstyle{same}
\hypersetup{
  pdftitle={几种低阶求积公式的余项},
  colorlinks=true,
  linkcolor={Maroon},
  filecolor={Maroon},
  citecolor={Blue},
  urlcolor={Blue},
  pdfcreator={LaTeX via pandoc}}

\title{几种低阶求积公式的余项}
\author{}
\date{}

\begin{document}
\maketitle

\subsubsection{4.2.3 几种低阶求积公式的余项}\label{ux51e0ux79cdux4f4eux9636ux6c42ux79efux516cux5f0fux7684ux4f59ux9879}

首先考察梯形公式，按余项公式（4.1.7），梯形公式（4.1.1）的余项为

\[
R_{\mathrm{T}}=I-T=\int_a^b\frac{f''(\xi)}2(x-a)(x-b)\,\mathrm{d}x.
\]

这里积分的核函数 \((x-a)(x-b)\) 在区间 \([a,b]\) 上保号（非正），应用积分中值定理，在 \((a,b)\) 内存在一点 \(\eta\)，使

\[
R_{\mathrm{T}}=\frac{f''(\eta)}2\int_a^b(x-a)(x-b)\,\mathrm{d}x
=-\frac{f''(\eta)}{12}(b-a)^3,\quad\eta\in(a,b).
\tag{4.2.5}
\]

再研究 Simpson 公式（4.2.3）的余项 \(R_{\mathrm{S}}=I-S\)。为此构造次数不大于三的多项式 \(H(x)\)，使之满足

\[
H(a)=f(a),\quad H(b)=f(b),\quad H(c)=f(c),\quad H'(c)=f'(c),
\tag{4.2.6}
\]

这里 \(c=\dfrac{a+b}{2}\)。由于 Simpson 公式具有三次代数精度，它对于这样构造出的三次式 \(H(x)\) 是准确的：

\[
\int_a^b H(x)\,\mathrm{d}x=\frac{b-a}{6}[H(a)+4H(c)+H(b)].
\]

而由插值条件（4.2.6）知，上式右端实际上等于按 Simpson 公式（4.2.3）求得的积分值 \(S\)，因此积分余项

\[
R_{\mathrm{S}}=I-S=\int_a^b[f(x)-H(x)]\,\mathrm{d}x.
\]

不难证明，对于满足条件（4.2.6）的多项式 \(H(x)\)，其插值余项

\[
f(x)-H(x)=\frac{f^{(4)}(\xi)}{4!}(x-a)(x-c)^2(x-b),
\]

故有

\[
R_{\mathrm{S}}=\int_a^b\frac{f^{(4)}(\xi)}{4!}(x-a)(x-c)^2(x-b)\,\mathrm{d}x.
\]

这时积分核 \((x-a)(x-c)^2(x-b)\) 在 \([a,b]\) 上保号（非正），用积分中值定理有

\[
R_{\mathrm{S}}=\frac{f^{(4)}(\eta)}{4!}\int_a^b(x-a)(x-c)^2(x-b)\,\mathrm{d}x
=-\frac{b-a}{180}\left(\frac{b-a}{2}\right)^4f^{(4)}(\eta).
\tag{4.2.7}
\]

关于 Cotes 公式（4.2.4）的积分余项，这里不再具体推导，仅列出结果如下：

\[
R_{\mathrm{C}}=I-C=-\frac{2(b-a)}{945}\left(\frac{b-a}{4}\right)^6f^{(6)}(\eta).
\tag{4.2.8}
\]

\end{document}
